\documentclass[11pt,a4paper]{article}

\usepackage[utf8]{inputenc}
\usepackage[T1]{fontenc}
\usepackage{lmodern}
\usepackage{amsmath,amssymb,amsthm,mathtools}
\usepackage{geometry}
\usepackage{hyperref}
\usepackage{cleveref}
\usepackage{enumitem}
\usepackage{xcolor}
\usepackage{booktabs}
\usepackage{longtable}
\usepackage{array}
\usepackage{fancyhdr}

\geometry{margin=2.5cm}

\newcommand{\Tr}{\operatorname{Tr}}
\newcommand{\tr}{\operatorname{tr}}
\newcommand{\dd}{\mathrm{d}}
\newcommand{\Id}{\mathbf{1}}
\newcommand{\Weylsq}{C_{\mu\nu\rho\sigma}C^{\mu\nu\rho\sigma}}
\newcommand{\Rsq}{R^2}
\newcommand{\Ricsq}{R_{\mu\nu}R^{\mu\nu}}
\newcommand{\Boxop}{\Box}
\DeclareMathOperator{\erfi}{erfi}

\newtheorem{theorem}{Theorem}[section]
\newtheorem{proposition}[theorem]{Proposition}
\newtheorem{lemma}[theorem]{Lemma}
\theoremstyle{definition}
\newtheorem{definition}[theorem]{Definition}
\theoremstyle{remark}
\newtheorem{remark}[theorem]{Remark}

\pagestyle{fancy}
\fancyhf{}
\fancyhead[L]{\small\textsc{SCT Theory --- Derivation Probe}}
\fancyhead[R]{\small\thepage}
\renewcommand{\headrulewidth}{0.4pt}

\hypersetup{
  colorlinks=true,
  linkcolor=blue!60!black,
  citecolor=green!40!black,
  urlcolor=blue!60!black,
  pdftitle={Local Scalar Sector in the Spectral Action},
  pdfauthor={David Alfyorov},
}

\title{%
  \textbf{Local Scalar Sector in the Spectral Action}\\[6pt]
  \large A Derivation Probe on the Weyl and $R^2$ Limits of the Scalar Contribution
}
\author{David Alfyorov}
\date{March 2026}

\begin{document}
\maketitle

\begin{center}
\small\textit{Credits: Aliaksandr Samatyia contributed research-assistance and workflow support.}
\end{center}

\thispagestyle{empty}

\begin{abstract}
This note is a derivation-style probe written to match the working rhythm of
the SCT theory corpus.  Its aim is narrow but structurally important: to
track, in full local detail, how a single real scalar field contributes to
the curvature-squared sector of the one-loop spectral action.  The two facts
to be established are familiar from the verified scalar form-factor phase,
but here they are reassembled in a way that reflects the style of the
derivation corpus itself:
\begin{enumerate}[nosep]
  \item the scalar contribution to the Weyl-squared sector has a fixed local
  coefficient $\beta_W^{(0)} = 1/120$;
  \item the scalar contribution to the local $R^2$ sector is
  $\beta_R^{(0)}(\xi) = \frac{1}{2}(\xi - 1/6)^2$ and therefore vanishes at
  conformal coupling.
\end{enumerate}
The derivation is carried out in the Weyl basis $\{\Weylsq,\Rsq\}$, beginning
from the verified scalar heat-kernel form factors and the master function
\[
  \phi(x) = \int_0^1 \dd\alpha\, e^{-\alpha(1-\alpha)x}
  = e^{-x/4}\sqrt{\frac{\pi}{x}}\,
    \erfi\!\left(\frac{\sqrt{x}}{2}\right).
\]
We keep a strict separation between exact formulas, local expansions,
structural interpretation, and domain of validity.  No new claim is made
beyond the already verified scalar sector; the point is to expose clearly how
the conformal locus enters the local spectral action.
\end{abstract}

\tableofcontents
\newpage

\section{Introduction}
\label{sec:intro}

The curvature-squared sector of the spectral action is one of the first
places where the formal machinery of the heat kernel begins to show direct
physical structure.  At that order, each field species contributes a pair of
form factors multiplying $\Weylsq$ and $\Rsq$, and the local limits of those
form factors become the corresponding Wilson coefficients in the effective
action.  For the scalar field this sector is especially instructive because
it is the only spin contribution that carries the non-minimal coupling
$\xi$.

That single fact has two consequences.

First, the scalar sector is the only place where the local $R^2$ coefficient
can inherit dependence on the Higgs-curvature coupling when one passes to the
full Standard Model sum.  Second, the conformal value $\xi = 1/6$ is not an
external ornament imposed on the theory after the fact.  It appears directly
in the local coefficient extracted from the scalar form factor.

The aim of the present note is therefore not computational novelty but
structural clarity.  We derive the local limits in a way that makes the
dependence on $\xi$ transparent and keeps the logic in the order used
throughout the stronger SCT derivations:
\begin{enumerate}[label=\arabic*.]
  \item fix conventions and operator data,
  \item state the verified form factors,
  \item expand the master function,
  \item compute the local coefficients explicitly,
  \item read off the limiting cases and their physical interpretation.
\end{enumerate}

\section{Setup and Notation}
\label{sec:setup}

\subsection{Scalar operator data}

We work on a closed four-dimensional Riemannian manifold $(M,g)$ and take a
real scalar field $\varphi$ with action
\begin{equation}
  S[\varphi]
  = \frac{1}{2}\int \dd^4x\,\sqrt{g}\,
  \left(
    g^{\mu\nu}\partial_\mu\varphi\,\partial_\nu\varphi
    + \xi R\,\varphi^2 + m^2\varphi^2
  \right).
  \label{eq:scalar-action}
\end{equation}
The associated Laplace-type operator is
\begin{equation}
  \Delta = -\Boxop + \xi R + m^2
  = -\bigl(g^{\mu\nu}\nabla_\mu\nabla_\nu + E\bigr),
  \qquad
  E = -\xi R - m^2.
  \label{eq:laplace-operator}
\end{equation}

At one loop, the curvature-squared sector of the effective action takes the
Weyl-basis form
\begin{equation}
  \Gamma^{(1)}_{\mathcal R^2}
  = \frac{1}{2}\int \dd^4x\,\sqrt{g}\,
  \left[
    h_C^{(0)}\!\left(\frac{\Boxop}{m^2}\right)\Weylsq
    + h_R^{(0)}\!\left(\frac{\Boxop}{m^2};\xi\right)\Rsq
  \right].
  \label{eq:effective-action}
\end{equation}
Everything in this note is about the local limit of the two scalar form
factors $h_C^{(0)}$ and $h_R^{(0)}$.

\subsection{Verified scalar form factors}

The scalar Weyl form factor is
\begin{equation}
  h_C^{(0)}(x)
  = \frac{1}{12x} + \frac{\phi(x)-1}{2x^2}.
  \label{eq:hC-scalar}
\end{equation}
The scalar $R^2$ form factor is
\begin{equation}
  h_R^{(0)}(x;\xi)
  = f_{R,\mathrm{bis}}(x)
  + \xi\,f_{RU}(x)
  + \xi^2 f_U(x),
  \label{eq:hR-scalar}
\end{equation}
with
\begin{align}
  f_R(x)
  &= \frac{\phi(x)}{32}
   + \frac{\phi(x)}{8x}
   - \frac{7}{48x}
   - \frac{\phi(x)-1}{8x^2},
  \label{eq:fR}\\[4pt]
  f_{RU}(x)
  &= -\frac{\phi(x)}{4}
   - \frac{\phi(x)-1}{2x},
  \label{eq:fRU}\\[4pt]
  f_U(x)
  &= \frac{\phi(x)}{2},
  \label{eq:fU}\\[4pt]
  f_{R,\mathrm{bis}}(x)
  &= f_R(x) + \frac{1}{6}f_{RU}(x) + \frac{1}{36}f_U(x).
  \label{eq:fRbis}
\end{align}

These expressions are taken here as verified input.  The present derivation
does not reconstruct them from the full Barvinsky--Vilkovisky heat-kernel
machinery.  It uses them to expose the local coefficient structure in a way
that is easy to read and hard to misstate.

\subsection{The master function}

All scalar form factors are built from the single master function
\begin{equation}
  \phi(x)
  = \int_0^1 \dd\alpha\, e^{-\alpha(1-\alpha)x}
  = e^{-x/4}\sqrt{\frac{\pi}{x}}\,
    \erfi\!\left(\frac{\sqrt{x}}{2}\right).
  \label{eq:phi-def}
\end{equation}
It is the local expansion of $\phi(x)$ that governs the apparent poles at
$x=0$.  Once that expansion is in hand, the entire local structure of the
scalar curvature-squared sector follows by direct substitution.

\section{Local Expansion of the Master Function}
\label{sec:master}

\begin{lemma}
The master function has the local expansion
\begin{equation}
  \phi(x)
  = 1 - \frac{x}{6} + \frac{x^2}{60}
    - \frac{x^3}{840} + \frac{x^4}{15120}
    + \mathcal O(x^5).
  \label{eq:phi-series}
\end{equation}
\end{lemma}

\begin{proof}
Expand the exponential in \eqref{eq:phi-def} termwise:
\begin{equation}
  e^{-\alpha(1-\alpha)x}
  = \sum_{n=0}^{\infty}
    \frac{[-\alpha(1-\alpha)x]^n}{n!}.
  \label{eq:expansion-exp}
\end{equation}
Then
\begin{equation}
  \phi(x)
  = \sum_{n=0}^{\infty}
    \frac{(-x)^n}{n!}
    \int_0^1 \dd\alpha\, \alpha^n(1-\alpha)^n.
  \label{eq:phi-beta}
\end{equation}
The integral is a beta function:
\begin{equation}
  \int_0^1 \dd\alpha\, \alpha^n(1-\alpha)^n
  = \frac{(n!)^2}{(2n+1)!}.
  \label{eq:beta-integral}
\end{equation}
Therefore
\begin{equation}
  \phi(x)
  = \sum_{n=0}^{\infty}
    (-1)^n\frac{n!}{(2n+1)!}x^n,
  \label{eq:phi-general-series}
\end{equation}
and the first few coefficients give \eqref{eq:phi-series}.
\end{proof}

\begin{remark}
The form \eqref{eq:phi-def} contains $\sqrt{x}$, but the Taylor series
\eqref{eq:phi-general-series} shows that the resulting function is perfectly
regular at the origin.  The square-root singularity is only apparent.
\end{remark}

\section{The Local Weyl Coefficient}
\label{sec:weyl}

\subsection{Cancellation of the apparent pole}

Substituting \eqref{eq:phi-series} into \eqref{eq:hC-scalar} gives
\begin{align}
  h_C^{(0)}(x)
  &= \frac{1}{12x}
   + \frac{1}{2x^2}
     \left(
       -\frac{x}{6} + \frac{x^2}{60}
       - \frac{x^3}{840} + \frac{x^4}{15120}
       + \mathcal O(x^5)
     \right)
  \notag\\[4pt]
  &= \frac{1}{12x}
   - \frac{1}{12x}
   + \frac{1}{120}
   - \frac{x}{1680}
   + \frac{x^2}{30240}
   + \mathcal O(x^3).
  \label{eq:hC-series}
\end{align}
The pole cancels exactly.  Therefore the scalar Weyl form factor admits a
finite local limit.

\subsection{Result}

\begin{proposition}
The scalar contribution to the local Weyl-squared coefficient is
\begin{equation}
  \boxed{\beta_W^{(0)} = h_C^{(0)}(0) = \frac{1}{120}.}
  \label{eq:betaW}
\end{equation}
\end{proposition}

\begin{proof}
It is the constant term of the expansion \eqref{eq:hC-series}.
\end{proof}

This coefficient is independent of $\xi$.  That independence is important.
The non-minimal coupling affects the scalar response to scalar curvature, but
it does not deform the local Weyl coefficient.  In the full Standard Model
sum, the $\xi$-dependence therefore enters only through the $R^2$ channel.

\section{The Local \texorpdfstring{$R^2$}{R2} Coefficient}
\label{sec:R2}

\subsection{Expansion of the component form factors}

We now expand the three scalar functions entering $h_R^{(0)}(x;\xi)$.

\paragraph{The function $f_R(x)$.}
Using \eqref{eq:phi-series},
\begin{align}
  \frac{\phi(x)}{32}
  &= \frac{1}{32}
   - \frac{x}{192}
   + \frac{x^2}{1920}
   - \frac{x^3}{26880}
   + \mathcal O(x^4),
  \label{eq:phi-over-32}\\[4pt]
  \frac{\phi(x)}{8x}
  &= \frac{1}{8x}
   - \frac{1}{48}
   + \frac{x}{480}
   - \frac{x^2}{6720}
   + \mathcal O(x^3),
  \label{eq:phi-over-8x}\\[4pt]
  -\frac{\phi(x)-1}{8x^2}
  &= \frac{1}{48x}
   - \frac{1}{480}
   + \frac{x}{6720}
   - \frac{x^2}{120960}
   + \mathcal O(x^3).
  \label{eq:phi-minus-1}
\end{align}
Combining these pieces with $-7/(48x)$,
\begin{align}
  f_R(x)
  &= \left(\frac{1}{8x} + \frac{1}{48x} - \frac{7}{48x}\right)
   + \left(\frac{1}{32} - \frac{1}{48} - \frac{1}{480}\right)
  \notag\\
  &\qquad
   + \left(-\frac{1}{192} + \frac{1}{480} + \frac{1}{6720}\right)x
   + \left(\frac{1}{1920} - \frac{1}{6720} - \frac{1}{120960}\right)x^2
   + \mathcal O(x^3)
  \notag\\[4pt]
  &= \frac{1}{120}
   - \frac{x}{336}
   + \frac{11x^2}{30240}
   + \mathcal O(x^3).
  \label{eq:fR-series}
\end{align}

\paragraph{The function $f_{RU}(x)$.}
From \eqref{eq:phi-series},
\begin{align}
  -\frac{\phi(x)}{4}
  &= -\frac{1}{4}
   + \frac{x}{24}
   - \frac{x^2}{240}
   + \frac{x^3}{3360}
   + \mathcal O(x^4),
  \label{eq:minus-phi-over-4}\\[4pt]
  -\frac{\phi(x)-1}{2x}
  &= \frac{1}{12}
   - \frac{x}{120}
   + \frac{x^2}{1680}
   - \frac{x^3}{30240}
   + \mathcal O(x^4).
  \label{eq:phi-1-over-2x}
\end{align}
Hence
\begin{equation}
  f_{RU}(x)
  = -\frac{1}{6}
  + \frac{x}{30}
  - \frac{x^2}{280}
  + \frac{x^3}{3780}
  + \mathcal O(x^4).
  \label{eq:fRU-series}
\end{equation}

\paragraph{The function $f_U(x)$.}
This one is immediate:
\begin{equation}
  f_U(x)
  = \frac{\phi(x)}{2}
  = \frac{1}{2}
  - \frac{x}{12}
  + \frac{x^2}{120}
  - \frac{x^3}{1680}
  + \mathcal O(x^4).
  \label{eq:fU-series}
\end{equation}

\paragraph{The function $f_{R,\mathrm{bis}}(x)$.}
Using \eqref{eq:fRbis} together with \eqref{eq:fR-series},
\eqref{eq:fRU-series}, and \eqref{eq:fU-series}, one finds
\begin{align}
  f_{R,\mathrm{bis}}(x)
  &= \left(\frac{1}{120}\right)
   + \frac{1}{6}\left(-\frac{1}{6}\right)
   + \frac{1}{36}\left(\frac{1}{2}\right)
   + \mathcal O(x)
  \notag\\
  &= -\frac{1}{180}
   + \frac{x}{3780}
   + \mathcal O(x^2).
  \label{eq:fRbis-series}
\end{align}
Thus the local value is
\begin{equation}
  f_{R,\mathrm{bis}}(0) = -\frac{1}{180}.
  \label{eq:fRbis-local}
\end{equation}

\subsection{Assembling the scalar \texorpdfstring{$R^2$}{R2} form factor}

Substituting the local values into \eqref{eq:hR-scalar} gives
\begin{equation}
  h_R^{(0)}(0;\xi)
  = -\frac{1}{180}
  - \frac{\xi}{6}
  + \frac{\xi^2}{2}.
  \label{eq:hR-local-raw}
\end{equation}
Completing the square yields
\begin{equation}
  h_R^{(0)}(0;\xi)
  = \frac{1}{2}
    \left(
      \xi^2 - \frac{\xi}{3} + \frac{1}{36}
    \right)
  = \frac{1}{2}\left(\xi - \frac{1}{6}\right)^2.
  \label{eq:hR-local-square}
\end{equation}

\begin{proposition}
The scalar contribution to the local $R^2$ coefficient is
\begin{equation}
  \boxed{
    \beta_R^{(0)}(\xi)
    = h_R^{(0)}(0;\xi)
    = \frac{1}{2}\left(\xi - \frac{1}{6}\right)^2.
  }
  \label{eq:betaR}
\end{equation}
\end{proposition}

\begin{remark}
Equation \eqref{eq:betaR} is stronger than the statement that the result is
quadratic in $\xi$.  It says that the local scalar curvature response is a
perfect square centered at conformal coupling.  That point is the unique
minimum of the local scalar $R^2$ coefficient.
\end{remark}

\section{Insertion into the Local Spectral Action}
\label{sec:spectral}

\subsection{Single real scalar}

For a single real scalar, the local curvature-squared contribution can now be
written directly as
\begin{equation}
  \Gamma^{(1)}_{\mathcal R^2,\mathrm{local}}
  = \frac{1}{2}\int \dd^4x\,\sqrt{g}\,
  \left[
    \frac{1}{120}\Weylsq
    + \frac{1}{2}\left(\xi - \frac{1}{6}\right)^2 \Rsq
  \right].
  \label{eq:local-scalar-action}
\end{equation}
This formula is often the most useful local summary of the scalar sector.  It
states, in one line, which piece is rigid and which piece is controlled by
the non-minimal coupling.

\subsection{Why the conformal point survives in the full theory}

In the full Standard Model sum, the local $R^2$ coefficient receives no
contribution from the Dirac or vector sectors.  The scalar sector is
therefore the entire carrier of the $\xi$-dependence.  For four real scalar
degrees of freedom, one obtains
\begin{equation}
  \alpha_R(\xi)
  = 4 \times \frac{1}{2}\left(\xi - \frac{1}{6}\right)^2
  = 2\left(\xi - \frac{1}{6}\right)^2.
  \label{eq:alphaR-full}
\end{equation}
This is precisely the structural reason that the conformal point remains
visible after the full Standard Model sum.  It is not washed out by the
fermionic and vector sectors because they carry zero local $R^2$ coefficient.

\section{Conversion to the \texorpdfstring{$\{R^2,\Ricsq\}$}{R2,Ricci2} Basis}
\label{sec:basis}

The Weyl basis is natural for the heat-kernel derivation, but many
weak-field discussions of quadratic gravity use the basis
$\{R^2,\Ricsq\}$.  The relation between the two is standard.  In four
dimensions,
\begin{equation}
  \Weylsq
  = E_4 + 2\Ricsq - \frac{2}{3}\Rsq,
  \label{eq:weyl-conversion}
\end{equation}
where
\begin{equation}
  E_4
  = R_{\mu\nu\rho\sigma}R^{\mu\nu\rho\sigma}
  - 4R_{\mu\nu}R^{\mu\nu}
  + R^2
  \label{eq:euler-density}
\end{equation}
is the Euler density.  Up to the topological term, the local quadratic action
\begin{equation}
  \Gamma_{\mathrm{loc}}
  = \int \dd^4x\,\sqrt{g}\,
  \left[
    \alpha_C\Weylsq + \alpha_R\Rsq
  \right]
  \label{eq:weyl-basis-local}
\end{equation}
is therefore equivalent to
\begin{equation}
  \Gamma_{\mathrm{loc}}
  \simeq
  \int \dd^4x\,\sqrt{g}\,
  \left[
    c_1 \Rsq + c_2 \Ricsq
  \right],
  \qquad
  c_1 = \alpha_R - \frac{2}{3}\alpha_C,
  \quad
  c_2 = 2\alpha_C.
  \label{eq:c1c2-conversion}
\end{equation}

\subsection{Single-scalar coefficients in the Ricci basis}

For one real scalar,
\begin{equation}
  \alpha_C^{(0)} = \frac{1}{120},
  \qquad
  \alpha_R^{(0)}(\xi) = \frac{1}{2}\left(\xi-\frac{1}{6}\right)^2.
  \label{eq:single-scalar-alphas}
\end{equation}
Hence
\begin{align}
  c_2^{(0)}
  &= 2\alpha_C^{(0)}
   = \frac{1}{60},
  \label{eq:c2-single}\\[4pt]
  c_1^{(0)}(\xi)
  &= \alpha_R^{(0)}(\xi) - \frac{2}{3}\alpha_C^{(0)}
   = \frac{1}{2}\left(\xi-\frac{1}{6}\right)^2 - \frac{1}{180}.
  \label{eq:c1-single}
\end{align}
The scalar propagating combination becomes
\begin{equation}
  3c_1^{(0)} + c_2^{(0)}
  = \frac{3}{2}\left(\xi-\frac{1}{6}\right)^2.
  \label{eq:single-scalar-prop}
\end{equation}
Thus even before the Standard Model sum, the scalar channel is already
organized around the conformal point.

\section{Limiting Cases}
\label{sec:limits}

\subsection{Minimal coupling}

At $\xi = 0$ one has
\begin{equation}
  \beta_R^{(0)}(0)
  = \frac{1}{2}\left(-\frac{1}{6}\right)^2
  = \frac{1}{72}.
  \label{eq:minimal}
\end{equation}
Thus minimal coupling produces a positive local $R^2$ term.

\subsection{Conformal coupling}

At $\xi = 1/6$,
\begin{equation}
  \beta_R^{(0)}\!\left(\frac{1}{6}\right) = 0,
  \qquad
  \beta_W^{(0)} = \frac{1}{120}.
  \label{eq:conformal}
\end{equation}
Hence the scalar leaves a local Weyl trace but no local $R^2$ trace.

\subsection{Large non-minimal coupling}

For $|\xi| \gg 1$,
\begin{equation}
  \beta_R^{(0)}(\xi)
  = \frac{\xi^2}{2} - \frac{\xi}{6} + \frac{1}{72}
  \sim \frac{\xi^2}{2}.
  \label{eq:large-xi}
\end{equation}
Therefore the local $R^2$ response is rapidly enhanced when the non-minimal
coupling is large.  In that regime the scalar curvature coupling is not a
small correction but the dominant driver of the local scalar channel.

\subsection{Flat-space reading}

When the background is flat, $R=0$ and $C_{\mu\nu\rho\sigma}=0$, so the local
curvature-squared action vanishes.  The point of the local coefficients is
therefore not that they survive in flat space, but that they control the
first response once curvature is turned on.

\section{Physical Interpretation and Domain of Validity}
\label{sec:interpretation}

The main physical statement of this note is modest but clean:
\begin{enumerate}[label=(\roman*)]
  \item the scalar contribution to the Weyl sector is rigid;
  \item the scalar contribution to the local $R^2$ sector is minimized at
  $\xi = 1/6$;
  \item all departure from that point is quadratic in the deviation from
  conformality.
\end{enumerate}

This should not be oversold.  The result concerns the local curvature-squared
sector only.  It does not by itself determine the full nonlocal propagator,
the entire-function structure, or the complete spectrum of gravitational
excitations.  Those questions require the corresponding nonlocal analysis.

What the present result does establish is more elementary and more secure:
the conformal point is already encoded in the verified scalar form factors at
the local level.  One does not need to invoke extra philosophical arguments
to see it; the square appears directly in the coefficient.

\section{Cross-Checks and Bridges to the Full Theory}
\label{sec:checks}

\subsection{Check 1: minimal and conformal values}

Two benchmark points are worth displaying separately because they recur
throughout the scalar-sector literature:
\begin{align}
  h_R^{(0)}(0;0)
  &= \frac{1}{2}\left(-\frac{1}{6}\right)^2
   = \frac{1}{72},
  \label{eq:check-minimal}\\[4pt]
  h_R^{(0)}\!\left(0;\frac{1}{6}\right)
  &= \frac{1}{2}\left(\frac{1}{6}-\frac{1}{6}\right)^2
   = 0.
  \label{eq:check-conformal}
\end{align}
The first benchmark confirms the standard minimal-coupling scalar result. The
second confirms that the conformal point removes the local $R^2$ term
completely.

\subsection{Check 2: positivity}

The coefficient \eqref{eq:betaR} is manifestly non-negative:
\begin{equation}
  \beta_R^{(0)}(\xi) \geq 0
  \qquad\text{for all } \xi \in \mathbb R.
  \label{eq:positivity}
\end{equation}
This does not by itself guarantee positivity of the full nonlocal scalar
channel, but it is the correct local sign check and one of the reasons the
conformal point is structurally distinguished.

\subsection{Check 3: one real scalar versus the Higgs multiplet}

The local derivation above is for one real scalar degree of freedom.  The
Higgs doublet contributes four real scalar degrees of freedom.  Therefore the
local Standard Model scalar contribution is
\begin{align}
  \alpha_{C,\mathrm{scalar}}
  &= 4\times \frac{1}{120}
   = \frac{1}{30},
  \label{eq:alphaC-scalar-sector}\\[4pt]
  \alpha_{R,\mathrm{scalar}}(\xi)
  &= 4\times \frac{1}{2}\left(\xi-\frac{1}{6}\right)^2
   = 2\left(\xi-\frac{1}{6}\right)^2.
  \label{eq:alphaR-scalar-sector}
\end{align}
Since the local Dirac and vector $R^2$ coefficients vanish, the full Standard
Model $R^2$ coefficient is exactly the same expression as the Higgs-sector
contribution.  That statement is one of the most transparent bridges between
the single-spin derivation and the combined Standard Model result.

\subsection{Check 4: relation to the scalar propagating sector}

In the local quadratic-gravity language, the scalar propagating channel is
controlled by the combination $3c_1+c_2$.  Once the Weyl-basis coefficients
are converted into the $\{R^2,\Ricsq\}$ basis, the scalar-sector dependence
enters through the same square:
\begin{equation}
  3c_1 + c_2
  \propto \left(\xi - \frac{1}{6}\right)^2.
  \label{eq:3c1c2-bridge}
\end{equation}
This is the local bridge from the present derivation note to the
prediction-style statement that the scalar mode decouples at conformal
coupling.

\subsection{Check 5: dimensional reading}

The local coefficients $\beta_W^{(0)}$ and $\beta_R^{(0)}$ are dimensionless,
as they must be in four dimensions when multiplying curvature-squared
invariants in the one-loop effective action.  This is a trivial check, but it
is also a useful guard against the most common kind of bookkeeping error in
curvature expansions: forgetting where the mass scale has been divided out.

\section{Summary of Results}
\label{sec:summary}

For a single real scalar field with non-minimal coupling $\xi$, the local
curvature-squared sector of the one-loop spectral action is controlled by
\begin{align}
  \beta_W^{(0)} &= \frac{1}{120},
  \label{eq:summary-betaW}\\[4pt]
  \beta_R^{(0)}(\xi)
  &= \frac{1}{2}\left(\xi - \frac{1}{6}\right)^2.
  \label{eq:summary-betaR}
\end{align}
The first coefficient is fixed and $\xi$-independent.  The second is a
positive semidefinite perfect square centered at conformal coupling.

This pair of facts is the complete local scalar story at quadratic order in
curvature.  It is small in scope, but it is exactly the piece of the theory
that later controls why the full Standard Model $R^2$ coefficient carries
only scalar-sector $\xi$-dependence.

\appendix

\section{Taylor Coefficients of the Master Function}
\label{app:taylor}

From \eqref{eq:phi-general-series},
\begin{equation}
  \phi(x)
  = \sum_{n=0}^{\infty}
    (-1)^n\frac{n!}{(2n+1)!}x^n.
  \label{eq:appendix-series}
\end{equation}
The first coefficients are listed in \cref{tab:taylor}.

\begin{table}[ht]
\centering
\caption{First coefficients in the local expansion of $\phi(x)$.}
\label{tab:taylor}
\begin{tabular}{cc}
\toprule
$n$ & coefficient of $x^n$ \\
\midrule
0 & $1$ \\
1 & $-1/6$ \\
2 & $1/60$ \\
3 & $-1/840$ \\
4 & $1/15120$ \\
5 & $-1/332640$ \\
\bottomrule
\end{tabular}
\end{table}

\section{Benchmark Table for the Local Scalar Sector}
\label{app:benchmarks}

\begin{table}[ht]
\centering
\caption{Local scalar-sector benchmarks collected in one place.}
\label{tab:benchmarks}
\begin{tabular}{lll}
\toprule
Quantity & Value & Role \\
\midrule
$h_C^{(0)}(0)$ & $1/120$ & local Weyl coefficient \\
$f_R(0)$ & $1/120$ & scalar CZ component \\
$f_{RU}(0)$ & $-1/6$ & scalar CZ component \\
$f_U(0)$ & $1/2$ & scalar CZ component \\
$f_{R,\mathrm{bis}}(0)$ & $-1/180$ & Weyl-basis scalar component \\
$h_R^{(0)}(0;\xi)$ & $\frac{1}{2}(\xi-1/6)^2$ & local $R^2$ coefficient \\
$h_R^{(0)}(0;0)$ & $1/72$ & minimal coupling check \\
$h_R^{(0)}(0;1/6)$ & $0$ & conformal coupling check \\
\bottomrule
\end{tabular}
\end{table}

\section{Detailed Algebra for the Local \texorpdfstring{$R^2$}{R2} Coefficient}
\label{app:detailed-R2}

For ease of later reuse, we collect the local algebra once more in compact
form.  Starting from \eqref{eq:fR-series}, \eqref{eq:fRU-series}, and
\eqref{eq:fU-series},
\begin{align}
  f_{R,\mathrm{bis}}(0)
  &= f_R(0) + \frac{1}{6}f_{RU}(0) + \frac{1}{36}f_U(0)
  \notag\\
  &= \frac{1}{120} + \frac{1}{6}\left(-\frac{1}{6}\right)
   + \frac{1}{36}\left(\frac{1}{2}\right)
  \notag\\
  &= \frac{1}{120} - \frac{1}{36} + \frac{1}{72}
   = -\frac{1}{180}.
  \label{eq:appendix-frbis}
\end{align}
Hence
\begin{align}
  h_R^{(0)}(0;\xi)
  &= -\frac{1}{180} - \frac{\xi}{6} + \frac{\xi^2}{2}
  \notag\\
  &= \frac{1}{2}
     \left(
       \xi^2 - \frac{\xi}{3} + \frac{1}{36}
     \right)
  \notag\\
  &= \frac{1}{2}\left(\xi-\frac{1}{6}\right)^2.
  \label{eq:appendix-complete-square}
\end{align}
The two benchmark values used repeatedly in the main text follow
immediately:
\begin{equation}
  h_R^{(0)}(0;0)=\frac{1}{72},
  \qquad
  h_R^{(0)}\!\left(0;\frac{1}{6}\right)=0.
  \label{eq:appendix-benchmarks}
\end{equation}

\section{Basis Conversion Formulae}
\label{app:basis}

For completeness, we record the standard quadratic-basis conversion used in
\cref{sec:basis}.  In four dimensions the Weyl tensor satisfies
\begin{equation}
  C_{\mu\nu\rho\sigma}
  = R_{\mu\nu\rho\sigma}
  - \frac{1}{2}\left(
      g_{\mu\rho}R_{\nu\sigma}
      - g_{\mu\sigma}R_{\nu\rho}
      - g_{\nu\rho}R_{\mu\sigma}
      + g_{\nu\sigma}R_{\mu\rho}
    \right)
  + \frac{R}{6}\left(
      g_{\mu\rho}g_{\nu\sigma}
      - g_{\mu\sigma}g_{\nu\rho}
    \right).
  \label{eq:weyl-def-app}
\end{equation}
Squaring and contracting gives
\begin{equation}
  \Weylsq
  = R_{\mu\nu\rho\sigma}R^{\mu\nu\rho\sigma}
  - 2\Ricsq + \frac{1}{3}\Rsq.
  \label{eq:weyl-square-app}
\end{equation}
Using the Euler density
\begin{equation}
  E_4
  = R_{\mu\nu\rho\sigma}R^{\mu\nu\rho\sigma}
  - 4\Ricsq + \Rsq,
  \label{eq:euler-app}
\end{equation}
one eliminates the Riemann-squared term and obtains
\begin{equation}
  \Weylsq = E_4 + 2\Ricsq - \frac{2}{3}\Rsq,
  \label{eq:weyl-to-ricci-app}
\end{equation}
which is the conversion used in \eqref{eq:c1c2-conversion}.

\begin{thebibliography}{99}

\bibitem{BarvinskyVilkovisky1987}
A.~O.~Barvinsky and G.~A.~Vilkovisky,
``Beyond the Schwinger--DeWitt technique: Converting loops into trees and
in-in currents,''
\emph{Nucl. Phys. B} \textbf{282} (1987) 163--188.

\bibitem{CodelloZanusso2013}
A.~Codello and O.~Zanusso,
``On the non-local heat kernel expansion,''
\emph{J. Math. Phys.} \textbf{54} (2013) 013513,
arXiv:\href{https://arxiv.org/abs/1203.2034}{1203.2034}.

\bibitem{Connes1994}
A.~Connes,
\emph{Noncommutative Geometry},
Academic Press, 1994.

\bibitem{ParkerToms2009}
L.~Parker and D.~Toms,
\emph{Quantum Field Theory in Curved Spacetime},
Cambridge University Press, 2009.

\bibitem{Vassilevich2003}
D.~V.~Vassilevich,
``Heat kernel expansion: User's manual,''
\emph{Phys. Rept.} \textbf{388} (2003) 279--360,
arXiv:\href{https://arxiv.org/abs/hep-th/0306138}{hep-th/0306138}.

\end{thebibliography}

\end{document}
